% year: תשפ"ג
% type: מבחן
% semester: א
% moed: א
% number: 1ב
% points: 10
% categories: func-limits
% source: exams/מבחנים/תשפג/מועד א תשפג.pdf

\begin{question}
לאילו ערכים של הפרמטר $a$ יש לפונקציה הבאה גבול ב-$x=0$?
\[ f(x)=\begin{cases} \dfrac{\sin(2x)}{x}, & x<0 \\[2mm] \dfrac{x+3}{4x+a}, & x\ge 0 \end{cases} \]
\end{question}

\begin{hint}
חשבו בנפרד את הגבולות החד-צדדיים ודרשו שיהיו שווים. שימו לב למקרה $a=0$.
\end{hint}

\begin{solution}
לפונקציה יש גבול ב-$x=0$ אם ורק אם שני הגבולות החד-צדדיים קיימים (סופיים) ושווים.

\textbf{גבול משמאל:} לפי הגבול היסודי $\lim_{t\to 0}\frac{\sin t}{t}=1$,
\[ \lim_{x\to 0^-}\frac{\sin(2x)}{x}=\lim_{x\to 0^-}2\cdot\frac{\sin(2x)}{2x}=2. \]

\textbf{גבול מימין:} אם $a\neq 0$, המונה והמכנה רציפים ב-$0$ והמכנה שונה מ-$0$ שם, ולכן
\[ \lim_{x\to 0^+}\frac{x+3}{4x+a}=\frac{3}{a}. \]
אם $a=0$, אז $\frac{x+3}{4x}\to +\infty$ כאשר $x\to 0^+$ (המונה שואף ל-$3$ והמכנה ל-$0^+$), והגבול מימין אינו סופי, לכן אין גבול.

לכן עבור $a\ne 0$ דרוש $\frac{3}{a}=2$, כלומר $a=\frac32$. (במקרה זה המכנה $4x+\frac32>0$ לכל $x\ge 0$, כך שהפונקציה מוגדרת היטב בסביבה ימנית של $0$.)

\textbf{תשובה:} לפונקציה יש גבול ב-$x=0$ רק עבור $a=\frac{3}{2}$, וערך הגבול הוא $2$.
\end{solution}
